\section{Contributions}
\label{sec:contributions}
\paragraph{\textbf{Alessio Demo (Badge number: 2142885)}}
\begin{itemize}[leftmargin=*]
    \item \textbf{First Deadline:}  Developed function database connection and instance creation of DuckDB database manager for each dataset using the ingest files (.duckdb files).
    Developed function for retrieving the SQL queries from .sql files, processing the queries and executing them on DuckDB database instance and finally store it in a json file. 
    \item \textbf{Second Deadline:} Implemented the Baseline NL function for interacting with the LLM with a backoff in case of gemini returns a free-plan limitation exception, save the results in FULL JSON format and next using it in the main script. Unit tests for the baseline\_tools script using pytest. RAG task implementation using external libraries for creating the embeddings and evaluate the similarity between them in relation to a given prompt.
    Developed scripts for getting the db schema from DuckDB database for each dataset and the SQL query parsing, that will be used in the GALOIS implementation in the next tasks.
    \item \textbf{ThirdDeadline:} Deployed the query parser that translates SQL queries to an abstract syntax tree that
    makes it more easy to manipulate. Developed the Galois class that implements the logic
    handling the choice of physical and logical optimization techniques, also implemented
    the class GaloisEstimator that interacts with LLM for confidence rates involved in the
    process. Furthermore builded up the post-processing process that handles the execution
    of original query on the data frame containing the LLM response, and JOIN operations
    between multiple tables. Added costs optimization logic that asks the LLM to provide
    data only for the attributes that are involved in each query, instead of asking the entire
    set of attributes that make up the table.
    Developed Experiment 4, Experiment 5 and Experiment 8.
\end{itemize}
Writing the reports sections reporting the complete pipeline followed for developing the stuff above.

\paragraph{\textbf{Francesco Pivotto (Badge number: 2158296)}}
\begin{itemize}[leftmargin=*]
    \item \textbf{First Deadline:} Implemented a modular project structure by relocating core logic to \texttt{src/} and consolidating configuration and database management logic within dedicated utility modules (\texttt{config/loaders.py}). 
    Established the framework for multi-provider support by adding initial integration for the \texttt{Grok LLM}, including dedicated configuration points (\texttt{.env}, \texttt{config.yaml}, \texttt{loaders.py}).
    
    \item \textbf{Second Deadline:} Completed initial implementation of the SQL baseline, including the core logic for connecting to \texttt{DuckDB} and structuring query execution via the Watsonx provider.
    Fixed critical issues in environment variable loading, ensuring robust and standardized retrieval of secrets and settings via the central configuration object (using \texttt{loaders.py}).
    Modularized the \texttt{sql\_baseline} module, incorporating the LCEL Pattern to prepare for dynamic LLM selection.
    Centralized general utility functions by moving helpers into \texttt{baseline\_tools}.
    Introduced the dedicated \textbf{\texttt{llm\_wrapper}} class to handle all LLM connection logic, API parameter management, and calls.

    \item \textbf{Third Deadline:} Implemented the core Galois Execution Engine ($\texttt{GaloisExecutor}$) and integrated its iterative physical operators ($\texttt{Table-Scan}$ and $\texttt{Key-Scan}$). \\
    Decoupled the entire schema management architecture into modular components \\ ($\texttt{base\_schema\_manager.py}$, $\texttt{schema\_factory.py}$).\\
    

\end{itemize}

\paragraph{\textbf{Giorgia Amato (Badge number: 2159999)}}
\begin{itemize}[leftmargin=*]
    \item \textbf{First Deadline:} Contributed to project initialization by setting up the repository structure, environment configuration and utility scripts. 
    Added EXPLAIN and EXPLAIN ANALYZE support with JSON serialization, implemented JSON generation for query outputs, and improved logging and test utilities.
     \item \textbf{Second Deadline:} Developed the Gemini-based SQL baseline, including IK/MC query handling, schema and raw-data injection into prompts, JSON parsing through Pydantic, and error/timeout management. 
     Refactored and extended unit tests for baseline tools and added new tests validating LLM interaction and Gemini SQL execution.
     \item \textbf{Third Deadline:} Implemented the iterative extension for both SQL and NL baselines, introducing configurable execution through a \texttt{max\_iter} parameter, deduplication across iterations, and early-stopping conditions when no new tuples are returned. 
     Standardized the aggregation of token usage and latency metrics for iterative runs and ensured full compatibility with the existing JSON output schema.
     Implemented the GALOIS-WO Table-Scan pipeline, including the GaloisExecutor, the iterative loop (first + iterative prompts), the stopping criteria, and the JSON parsing/deduplication logic ensuring unique tuple collection. 
     Developed the GaloisWOSchemaManager for DuckDB schema inspection, table/attribute discovery, table-name resolution, key extraction, and JSON schema example generation. 
     Contributed to the writing and refinement of the report sections describing the iterative baselines, the GALOIS-WO Table-Scan algorithm, and the corresponding numerical results.
\end{itemize}